\documentclass[tikz,border=8pt]{standalone}
\usepackage{fontspec}
\usepackage{xeCJK}
\IfFontExistsTF{Arial}{\setmainfont{Arial}}{\setmainfont{Latin Modern Sans}}
\IfFontExistsTF{Microsoft JhengHei}{\setCJKmainfont{Microsoft JhengHei}}{\IfFontExistsTF{Noto Sans CJK TC}{\setCJKmainfont{Noto Sans CJK TC}}{\setCJKmainfont{SimSun}}}
\usetikzlibrary{arrows.meta,positioning,calc,fit,backgrounds,shapes.geometric}
\definecolor{ink}{HTML}{172033}
\definecolor{blue}{HTML}{2563EB}
\definecolor{gold}{HTML}{C79A2B}
\definecolor{green}{HTML}{16A34A}
\tikzset{
  card/.style={rounded corners=4pt, draw=ink!30, line width=.7pt, fill=white, minimum height=1.05cm, align=center, inner sep=6pt, text=ink},
  smallcard/.style={rounded corners=4pt, draw=ink!30, line width=.65pt, fill=white, minimum width=2.35cm, minimum height=.95cm, align=center, inner sep=5pt, text=ink, font=\small},
  goodarrow/.style={-Latex, draw=blue!75!black, line width=1.0pt},
  badarrow/.style={-Latex, draw=red!70!black, line width=1.0pt},
  soft/.style={draw=ink!35, line width=.8pt},
}
\begin{document}
\begin{tikzpicture}[font=\sffamily, >=Latex, line cap=round, line join=round]
\node[font=\bfseries\Large,text=ink] at (0,4.45) {Vibe Coding：從中文需求到會計洞察};
\node[card,fill=gold!18,draw=gold!70!black,minimum width=2.8cm] (n) at (-5.1,2.25) {中文需求};
\node[card,fill=blue!8,draw=blue!40,minimum width=2.8cm] (p) at (-2.55,2.25) {AI 產生\\程式草稿};
\node[card,fill=blue!8,draw=blue!40,minimum width=2.8cm] (t) at (0,2.25) {執行測試\\修錯};
\node[card,fill=green!8,draw=green!45!black,minimum width=2.8cm] (data) at (2.55,2.25) {資料表\\圖表};
\node[card,fill=green!8,draw=green!45!black,minimum width=2.8cm] (insight) at (5.1,2.25) {會計解讀};
\foreach \a/\b in {n/p,p/t,t/data,data/insight}{\draw[goodarrow] (\a) -- (\b);}
\node[smallcard,fill=red!6,draw=red!35,minimum width=3.7cm] at (-3.7,-0.15) {不要只交程式\\要交錯誤修正紀錄};
\node[smallcard,fill=red!6,draw=red!35,minimum width=3.7cm] at (0,-0.15) {不要只交資料\\要檢查來源與定義};
\node[smallcard,fill=red!6,draw=red!35,minimum width=3.7cm] at (3.7,-0.15) {不要只畫圖\\要回答會計問題};
\node[card,fill=white,draw=ink!25,minimum width=8.6cm] at (0,-2.45) {程式能力的入口，可以是清楚描述問題的能力};
\end{tikzpicture}
\end{document}
